\chapter*{ТЕРМИНЫ И ОПРЕДЕЛЕНИЯ}
\addcontentsline{toc}{chapter}{ТЕРМИНЫ И ОПРЕДЕЛЕНИЯ}

\noindent Анализ тональности текста – подкласс методов обработки естественного языка для автоматической классификации текста по категориям позитивного, негативного и нейтрального характера

\noindent Браузерный инференс – выполнение вывода нейросетевой модели непосредственно в браузере на устройстве пользователя без обращения к серверной инфраструктуре

\noindent Дистилляция знаний – метод компрессии, при котором компактная модель-«ученик» обучается воспроизводить поведение крупной модели-«учителя», усваивая распределения её выходного слоя

\noindent Дообучение – адаптация предобученной модели к конкретной задаче посредством дополнительного обучения на целевом размеченном корпусе с обновлением весовых параметров

\noindent Квантование – метод оптимизации, преобразующий параметры модели в форматы пониженной разрядности (INT8, INT4) для сокращения объёма памяти и ускорения вычислений

\noindent Мультиязычная модель – языковая модель, обученная на корпусах нескольких языков одновременно и формирующая единое векторное представление для типологически разнородных языков

\noindent Токенизация – разбиение текста на элементарные единицы (токены) для последующего преобразования в числовые идентификаторы, подаваемые на вход модели

\noindent Трансформер – архитектура нейронной сети, основанная на механизме самовнимания, вычисляющем зависимости между токенами за один параллельный проход без рекуррентных соединений

\noindent WebAssembly – бинарный формат инструкций для виртуальной машины браузера, обеспечивающий выполнение скомпилированного кода с производительностью, близкой к нативной